\documentclass[11pt,a4paper]{article}
\usepackage[top=2.5cm,bottom=2.8cm,left=2.5cm,right=2.5cm]{geometry}
\usepackage{amsmath,amssymb,amsthm}
\usepackage{url}
\usepackage{hyperref,xcolor,enumitem,booktabs,array}

\hypersetup{colorlinks=true,linkcolor={blue!60!black},
  citecolor={blue!60!black},urlcolor={blue!60!black}}

\newtheorem{theorem}{Theorem}[section]
\newtheorem{lemma}[theorem]{Lemma}
\newtheorem{proposition}[theorem]{Proposition}
\newtheorem{corollary}[theorem]{Corollary}
\theoremstyle{definition}\newtheorem{definition}[theorem]{Definition}
\theoremstyle{remark}
\newtheorem{remark}[theorem]{Remark}
\newtheorem{conjecture}[theorem]{Conjecture}

\newcommand{\vp}{\varphi}
\newcommand{\lam}{\lambda}
\newcommand{\TCMB}{T_{\mathrm{CMB}}}
\newcommand{\Lcond}{\Lambda_{\mathrm{cond}}}
\newcommand{\SU}{\mathrm{SU}}
\newcommand{\SO}{\mathrm{SO}}
\newcommand{\twoI}{2I}
\newcommand{\twoT}{2T}
\newcommand{\twoO}{2O}
\newcommand{\twoD}{2D}
\newcommand{\QH}{Q_H}
\newcommand{\Qgr}{Q_{\mathrm{group}}}
\newcommand{\Ztwo}{\mathbb{Z}_2}
\newcommand{\Zthree}{\mathbb{Z}_3}
\newcommand{\Zfour}{\mathbb{Z}_4}
\newcommand{\cdir}{c_{\mathrm{dir}}}
\newcommand{\Odm}{\Omega_{\mathrm{DM}}}
\newcommand{\Ob}{\Omega_{b}}

\title{\textbf{The Beyond--Standard-Model D-type Sector of the
Density-Feedback Hopfion:\\ A Sterile, Dirac Dark-Matter Candidate}}
\author{Frederick Manfredi}
\date{2026 --- DRAFT}

\begin{document}
\maketitle

\begin{abstract}
The matter sectors of the density-feedback Hopfion are the finite
subgroups of the binary icosahedral vacuum group $\twoI$: the neutrino
($\QH=1$, cyclic $\Ztwo$, A-type), the charged lepton ($\QH=2$, $\twoI$
itself, $\widehat{E}_8$), and the quark/baryon ($\QH=3$, binary
tetrahedral $\twoT\subset\twoI$ at pentagon index $5$, $\widehat{E}_6$).
We make two structural observations and then develop, honestly and
conditionally, their dark-matter consequence. \emph{(i)~ADE completion.}
Reading the Standard-Model matter content as subgroups of $\twoI$ forces
a definite exclusion at the top of the ADE tier: the binary octahedral
group $\twoO$ ($\widehat{E}_7$) cannot carry a matter sector, since
$|\twoO|=48$ does not divide $|\twoI|=120$. The framework predicts no
$E_7$-type matter. The binary dihedral (D-type) groups, by contrast, are
subgroups of $\twoI$ that carry their own $\Ztwo$ rather than the colour
$\Zthree$, and are the natural home of any beyond--Standard-Model state.
\emph{(ii)~A sterile Dirac sector.} The D-type candidates $\twoD_3$ and
$\twoD_5$ (McKay algebras $\SO(10)$ and $\SO(14)$) fail the T-matrix
closure that defines a quantised sector charge $\Qgr$, by both natural
WZW assignments; the obstruction is the odd-rank spinor, which also makes
the spinor \emph{complex}. The same object therefore renders the sector
both \emph{sterile} (no quantised Standard-Model charge) and
\emph{Dirac} (a conserved dark-number $U(1)$). \emph{(iii)~Dark matter,
as a developed conditional scenario.} Such a sector supplies a
gravitating, electromagnetically dark, cold ($a^{-3}$) relic, with an
abundance that an asymmetric-dark-matter co-genesis---one chiral,
CP-violating reconnection in the shared $\twoI$ vacuum seeding correlated
baryon and dark asymmetries---would set to $\Odm/\Ob\approx5$, i.e.\
$m_{\mathrm{dark}}\approx5$~GeV. We state precisely the four conditions
this scenario requires and which remain unestablished (the baryogenesis
magnitude $\eta$; the dark/baryon branching; the annihilation efficiency;
the soliton's \emph{existence}---its \emph{stability}, by contrast, is
topological against unwinding and barrier-protected (metastable) against
decay, absolute only if its $\Ztwo$ parity is dynamical). This paper is a
draft: results
(i)--(ii) are firm; the dark-matter application (iii) is a structured,
falsifiable, but conditional proposal, not a derivation of $\Odm$.
\end{abstract}

\tableofcontents

%══════════════════════════════════════════════════════════════════════
\section{Introduction and epistemic scope}
%══════════════════════════════════════════════════════════════════════
The density-feedback Hopfion~\cite{Paper1} organises the Standard-Model
matter content as topological excitations of a binary-icosahedral vacuum
condensate; Papers~XV--XX~\cite{Paper15,Paper16,Paper17,Paper18,Paper19,Paper20}
develop its three matter sectors. These are identified with
three finite subgroups of the binary icosahedral group $\twoI$ via the
McKay correspondence~\cite{CIZ1987}:
$\QH=1$ (neutrino, $\Ztwo$, $\widehat{A}_1$, $\SU(2)_1$), $\QH=2$
(charged lepton, $\twoI$, $\widehat{E}_8$, $\SU(2)_3$), and $\QH=3$
(quark/baryon, $\twoT$, $\widehat{E}_6$, $(E_6)_1\supset\SU(3)_1^{\times3}$).
The companion gravity paper~VII~\cite{Paper7} reproduces the galactic
dark-matter phenomenology (flat rotation curves, the radial-acceleration
relation) from director strain, but leaves the recombination-epoch
cold-dark-matter budget as an input. This paper asks whether the
framework's own group-theoretic ladder contains a candidate for that
budget.

\paragraph{Epistemic scope (stated up front).} We separate firm results
from a conditional proposal. Sections~\ref{P21:sec:ade}--\ref{P21:sec:dirac}
(the ADE completion, the no-$E_7$ exclusion, and the sterile-Dirac
character of the D-type sector) are firm: they are group theory and
rational conformal-field-theory arithmetic. Sections~\ref{P21:sec:dm}--%
\ref{P21:sec:abundance} (the dark-matter interpretation and its abundance)
are a \emph{developed conditional scenario}: every link is
framework-native, but the chain rests on quantities not yet derived,
which we enumerate explicitly in Section~\ref{P21:sec:open}. We flag the
status of each claim in place.

%══════════════════════════════════════════════════════════════════════
\section{The ADE completion and the exclusion of $E_7$}
\label{P21:sec:ade}
%══════════════════════════════════════════════════════════════════════
The matter sectors are not free McKay labels: they are \emph{subgroups}
of the vacuum group $\twoI$. Paper~XV~\cite{Paper15}, Theorem~9.1
%\ref{P15:thm:subgroup}
proves $\twoT\subset\twoI$ (orders $24\mid120$), with index
$[\twoI:\twoT]=5$ equal to the pentagon $k+2$ at $k=3$, the five cosets
being the five lepton primaries; the restriction $\twoI\to\twoT$ is the
algebraic content of passing from the lepton to the quark sector.
The neutrino group $\Ztwo=\{\pm1\}$ is the centre of $\SU(2)$ and lies
in every binary polyhedral group, $\Ztwo\subset\twoI$. Thus all three
matter sectors are $\twoI$ or subgroups of it.

\begin{theorem}[No $E_7$-type matter sector]
\label{P21:thm:noE7}
If matter sectors are realised as finite subgroups of the vacuum group
$\twoI$, then the binary octahedral group $\twoO$ (McKay $\widehat{E}_7$)
carries no matter sector.
\end{theorem}
\begin{proof}
$|\twoO|=48$ and $|\twoI|=120$. Since $48\nmid120$, Lagrange's theorem
forbids $\twoO\subset\twoI$. Equivalently, the octahedral group
$O\cong S_4$ (order $24$) is not a subgroup of the icosahedral group
$I\cong A_5$ (order $60$, with $24\nmid60$), so no lift to the binary
cover exists. Hence no $\widehat{E}_7$ matter sector is available.
\end{proof}

\noindent This is a definite, falsifiable prediction, independent of the
dark-matter discussion below: the Hopfion ladder admits A-type, $E_6$,
and $E_8$ matter, but never $E_7$.

\begin{remark}[The D-type slot is available and colourless]
\label{P21:rem:dtype_slot}
The binary dihedral (dicyclic) groups $\twoD_n$ of order $4n$ are
subgroups of $\twoI$ whenever $4n\mid120$ after the appropriate lift;
in particular $\twoD_3$ (order $12$, $\widehat{D}_5$) and $\twoD_5$
(order $20$, $\widehat{D}_7$) lie in $\twoI$ but \emph{not} in $\twoT$
(since $A_4$ has no subgroup of order $6$, neither dicyclic lifts into
$\twoT$). They are ``sideways branches'' of the vacuum. Crucially,
D-type groups carry their own $\Ztwo$ outer automorphism and not the
colour $\Zthree$ (Paper~XVII~\cite{Paper17}, A-type versus E-type McKay
chains): they are colourless, non-baryonic, and
reserved for beyond--Standard-Model physics. We take $\twoD_3,\twoD_5$
as the candidate dark sectors.
\end{remark}

%══════════════════════════════════════════════════════════════════════
\section{The D-type sector is sterile}
\label{P21:sec:sterile}
%══════════════════════════════════════════════════════════════════════
A sector carries a quantised charge iff its WZW T-matrix \emph{closes}:
following Paper~XV~\cite{Paper15} (Definition of T-matrix closure), an
integer $Q\geq1$ closes the T-matrix if $Q\,T_j$ is an odd multiple of
$\tfrac14$ for every primary $j$, where $T_j=h_j-c/24$; the minimal such
$Q$ is $\Qgr$. The Standard-Model sectors close: $(E_6)_1\to\Qgr=3$,
$\SU(2)_1\to\Qgr=6$, $\SU(2)_3\to\Qgr=10$ (reproduced by the engine of
Appendix~\ref{P21:app:engine}).

\begin{proposition}[No closing charge for the D-type sector]
\label{P21:prop:noclosure}
Neither natural WZW assignment of the D-type sector admits a closing
integer:
\begin{enumerate}[noitemsep,label=(\alph*)]
\item $(D_{n+2})_1=\SO(2r)_1$ (the McKay-algebra convention, as the quark
uses $(E_6)_1$): for $\twoD_3\to\SO(10)_1$ ($r=5$) and
$\twoD_5\to\SO(14)_1$ ($r=7$), the spinor primary has $h_s=r/8$, forcing
$Q\,r$ odd, while the vacuum $T_1=-r/24$ forces $Q\,r$ even. No integer
$Q$ satisfies both.
\item $\SU(2)_k$ with the D-type modular invariant (the convention the
lepton uses, $k=2n$): all even levels $k$ fail closure, while only odd
$k$ close; the D-type invariants sit at $k=6$ ($\twoD_3$) and $k=10$
($\twoD_5$), both even.
\end{enumerate}
Hence $\Qgr$ does not exist for the D-type sector.
\end{proposition}
\begin{proof}
Direct computation of the closure condition on the SO$(2r)_1$ primaries
$\{1,v,s,c\}$ with $(h_1,h_v,h_s)=(0,\tfrac12,\tfrac r8)$ and on the
$\SU(2)_k$ primaries $h_\ell=\ell(\ell+2)/4(k+2)$; see
Appendix~\ref{P21:app:engine}, which validates against the three
Standard-Model values before reporting the null results.
\end{proof}

\begin{corollary}[Sterility]
\label{P21:cor:sterile}
T-matrix closure is the charge-quantisation condition; its failure means
the D-type sector carries \emph{no} quantised Standard-Model charge.
Equivalently, electric charge is the $\SU(3)_L$ generator
$Q=T_{3L}+T_{8L}/\sqrt3$ of the $E_6$ trinification
(Paper~XV~\cite{Paper15}), quantised in units of $\tfrac13$ by the
$\Zthree$ centre of $\SU(3)_L$; a D-type group carries only its own
$\Ztwo$ and no $\Zthree$, so it admits no $\SU(3)_L$ representation for
$Q$ to act on. Being a separate McKay branch, it has no Standard-Model
doublet partner and hence no weak charge either. The D-type sector is a
full Standard-Model singlet: \emph{sterile}, interacting only
gravitationally and through its own $\Ztwo$.
\end{corollary}

%══════════════════════════════════════════════════════════════════════
\section{The same spinor makes the sector Dirac}
\label{P21:sec:dirac}
%══════════════════════════════════════════════════════════════════════
The obstruction in Proposition~\ref{P21:prop:noclosure}(a) is the odd-rank
spinor $h_s=r/8$. The reality of SO$(2r)$ spinors is classical: for $r$
odd the two Weyl spinors are \emph{complex} conjugates of one another
($s^{*}=c\neq s$); for $r\equiv0\ (\mathrm{mod}\ 4)$ real, for
$r\equiv2\ (\mathrm{mod}\ 4)$ pseudoreal. Both candidates have $r$ odd
($\SO(10):r=5$, $\SO(14):r=7$), so both carry a \emph{complex} spinor.
This is the same fact that makes the $\SO(10)$ \textbf{16} complex and
chiral, i.e.\ that makes $\SO(10)$ a viable grand-unified group.

\begin{proposition}[Sterile and Dirac share one origin]
\label{P21:prop:sterile_dirac}
The odd-rank spinor $h_s=r/8$ is simultaneously (i) the primary whose
T-matrix phase obstructs closure (Proposition~\ref{P21:prop:noclosure}a),
making the sector sterile, and (ii) complex ($s^{*}=c$), endowing the
sector with a conserved dark-number $U(1)$ and hence a distinct
antiparticle. A complex (Dirac) sector can carry a particle--antiparticle
asymmetry; contrast the neutrino sector $\SU(2)_1$, whose $j=\tfrac12$
primary is self-conjugate ($j^{*}=k-j=\tfrac12$), i.e.\ Majorana
(Paper~XVII~\cite{Paper17}), and which therefore cannot.
\end{proposition}

\noindent The D-type sector is thus a \emph{sterile Dirac} sector carrying
its own $\Ztwo$ (a dark parity, stabilising its lightest state) and a
complex-spinor $U(1)$ (a dark number, admitting an asymmetry). Its McKay
algebra for $\twoD_3$ is $\SO(10)$---the archetypal grand-unified group,
the home of baryon-number violation---a point we return to in
Section~\ref{P21:sec:abundance}. (We use $\SO(10)$ here as the
McKay-correspondent algebra, not as a gauged grand-unified theory.)

%══════════════════════════════════════════════════════════════════════
\section{Dark-matter candidacy: coldness and equation of state}
\label{P21:sec:dm}
%══════════════════════════════════════════════════════════════════════
A sterile, gravitating sector is a dark-matter candidate if it can be
\emph{cold} at recombination, i.e.\ dilute as $a^{-3}$. The D-type mass
scale is not fixed by the tower (there is no $\Qgr$ to set it), so we do
not predict it here; instead we identify the condition under which the
sector's gravitating content is $a^{-3}$.

\begin{proposition}[$a^{-3}$ equation of state, conditional]
\label{P21:prop:eos}
Let the sector's gravitating content be a population of localised,
number-conserved coherent configurations of fixed proper size and energy
(solitons, or frozen director deformations held rigid by the condensate's
static shear rigidity, Paper~VII~\cite{Paper7}). Then the number density
dilutes as $a^{-3}$ and the energy density as $\rho\propto a^{-3}$,
$w=0$: cold dark matter. If instead the content relaxes to a smooth
space-filling texture, gradients dilute as $a^{-1}$ and $\rho\propto
a^{-2}$, $w=-\tfrac13$: not cold.
\end{proposition}

\noindent Which case obtains is set by whether the configurations stay
localised or coarsen. In the framework's inertial director dynamics
($\chi\,\partial_t^2\hat n=K\nabla^2\hat n$), coarsening proceeds at a
rate controlled by the signal speed $\cdir=\sqrt{K/\chi}$, with no true
static plateau but a coarsening time $t_{\mathrm{coarsen}}\sim
\cdir^{-3/2\text{--}2}$ that diverges as $\cdir\to0$
(Paper~VII~\cite{Paper7}; the $3$D inertial defect-network simulation).
A coarse single-quench freeze test confirms the localised ($a^{-3}$)
regime for small $\cdir$ and coarsening for larger $\cdir$, with a sharp
threshold; repeated mergers (re-quenching) re-seed and \emph{sustain} the
localised network. The $a^{-3}$ regime therefore holds at recombination
provided $\cdir$ is below the sector's coldness threshold---the same
coldness the director sector already requires. This is a genuine escape
from the $a^{-2}$ obstruction that defeats a smooth orientation texture,
and it is purchased with no new ingredient.

%══════════════════════════════════════════════════════════════════════
\section{Abundance: an asymmetric-dark-matter co-genesis}
\label{P21:sec:abundance}
%══════════════════════════════════════════════════════════════════════
We summarise the abundance scenario and its conditions here; the genesis
dynamics---baryogenesis and co-genesis at the first-order confinement
transition---are developed in the companion Paper~XXII~\cite{Paper22},
which serves as the bridge from the quark sector to this dark sector. The
observed coincidence is $\Odm/\Ob\approx5.4$.

\paragraph{Baryogenesis as chiral reconnection.} In the framework a
baryon is a chiral wound state: the trefoil $T_{2,3}$ (writhe $+3$)
differs from its mirror (writhe $-3$), and is assembled by a chiral
threading process (Paper~XVIII~\cite{Paper18}). The three Sakharov
conditions map onto knot dynamics: (i)~baryon-number violation is the
reconnection $\QH{=}2+\QH{=}2\to\QH{=}3+\QH{=}1$ and its reverse (baryons
form from leptons by threading; $B$ is a knot-sector count changed by
reconnection, Paper~XVIII); (ii)~C and CP violation are supplied by the
derived WZW phase $\delta_{\mathrm{CP}}$ (Paper~V~\cite{Paper5}) together
with the chiral writhe, which bias the threading handedness; (iii)~the
departure from equilibrium is the confinement transition, which freezes
the reversible reconnections (``confinement as threading
irreversibility,'' Paper~XVIII). Because the $B$-violation here is
topological reconnection at the confinement scale and not an electroweak
sphaleron, the mechanism is not bound by the standard objection that the
Cabibbo--Kobayashi--Maskawa phase is too small for electroweak
baryogenesis.

\paragraph{The connector.} The phase $\delta_{\mathrm{CP}}$ is a property
of the $\twoI$ WZW T-matrix, common to the whole vacuum. The same chiral,
CP-violating reconnection event therefore biases both the $\twoT$ (baryon)
and the $\twoD_3$ (dark) freeze-outs that occur in the shared vacuum,
seeding \emph{correlated} asymmetries---a co-genesis. The D-type sector
can carry the dark asymmetry because it is Dirac
(Proposition~\ref{P21:prop:sterile_dirac}).

\paragraph{Overproduction is a feature.} A Kibble--Zurek freeze-out
produces a symmetric dark population of density $\sim n_\gamma$, far
exceeding the baryon asymmetry. In an asymmetric-dark-matter scenario
this is expected: the symmetric part self-annihilates---the D-type sector
is self-interacting through the shared director field---leaving the
CP-biased asymmetric residue. If that residue tracks the baryon asymmetry
(same $\delta_{\mathrm{CP}}$, comparable subgroup branching
$\twoD_3:\twoT$ of orders $12:24$), then
\begin{equation}
  \frac{\Odm}{\Ob}=r\,\frac{m_{\mathrm{dark}}}{m_p},\qquad
  r=\frac{n_{\mathrm{dark}}}{n_b},
  \label{P21:eq:adm}
\end{equation}
where the gravitating mass per relic is its rest mass $m_{\mathrm{dark}}$:
the relic is a stable knot (Hopf-charge conservation forbids its
unwinding), so no fraction of its energy is lost to a frozen residue.
$\Odm/\Ob\approx5.4$ with $r\approx1$ then gives
$m_{\mathrm{dark}}\approx5\,m_p\approx5$~GeV---the canonical
asymmetric-dark-matter mass. A $\sim5$~GeV sterile state evades the
null results of direct detection by construction. Non-thermal
(topological) production is consistent with $m_{\mathrm{dark}}\gg T_c$,
exactly as the $\sim1$~GeV proton forms at $T_c\sim150$~MeV.

%══════════════════════════════════════════════════════════════════════
\section{Falsifiable signatures}
\label{P21:sec:signatures}
%══════════════════════════════════════════════════════════════════════
\begin{enumerate}[noitemsep,label=(\arabic*)]
\item \textbf{No $E_7$ matter} (Theorem~\ref{P21:thm:noE7})---firm, and
independent of the dark-matter interpretation.
\item A $\sim5$~GeV sterile, asymmetric dark-matter state (if the
abundance scenario holds)---accessible in principle to light-dark-matter
searches, though sterile and so weakly coupled.
\item \textbf{Self-interacting dark matter}: the shared director field
makes dark--dark interactions unavoidable, predicting cored (not cuspy)
halos and the standard self-interacting-dark-matter small-scale
signatures.
\item \textbf{Merger dependence}: because the frozen network slowly
coarsens and mergers re-seed it, recently-merged/disturbed systems carry
a denser network, hence a larger mass excess, than relaxed systems of
equal baryon content (shared with Paper~VII~\cite{Paper7}'s cluster
channel).
\end{enumerate}

%══════════════════════════════════════════════════════════════════════
\section{Open problems}
\label{P21:sec:open}
%══════════════════════════════════════════════════════════════════════
The dark-matter application of Sections~\ref{P21:sec:dm}--\ref{P21:sec:abundance}
is a developed conditional scenario. Four quantities remain unestablished,
in decreasing order of difficulty.
\begin{enumerate}[noitemsep,label=(\roman*)]
\item \textbf{The baryogenesis magnitude $\eta$.} The Sakharov structure
is in place but $\eta\approx6\times10^{-10}$ is not derived. Paper~XXII%
~\cite{Paper22} develops this: the CP interference and the departure from
equilibrium are both $O(1)$, so the efficient reconnection
\emph{overproduces}, and the required suppression is carried by the
heavy-band anisotropy $S_{\mathrm{eff}}\propto\sin^4\theta/\vp^6$ of the
baryon channel, set by the same density-feedback strength $\beta_*\rho$
that drives the first-order transition. A self-consistent value inside the
first-order window reproduces $\eta$ but is not forced; what remains open
is the window's location in physical units and the heavy-band power. This
is the single lynchpin: with it the scenario becomes a prediction, without
it it stays conditional.
\item \textbf{The soliton's existence} (not its stability). The $a^{-3}$
and asymmetric-residue pictures require a localised D-type configuration
that is an energy minimiser in its sector. This \emph{existence} question
is open: coarse free-flow relaxation cannot settle it, because it does
not conserve the Hopf charge on a lattice---it fails even for the known
trefoil---so a charge-conserving relaxation is required. \emph{Stability},
a separate matter, splits into two questions of different strength. Against
\emph{unwinding} ($\QH\to0$) the configuration is protected topologically
by Hopf-charge conservation, exactly as the trefoil is; the charge loss
seen in free-flow relaxation is a lattice artefact of the non-conserving
scheme, not a dynamical unwinding. Against \emph{decay into
Standard-Model sectors} the firm argument is geometric: $\twoD_3$ and
$\twoD_5$ are sideways branches---subgroups of $\twoI$ but not of $\twoT$
(Remark~\ref{P21:rem:dtype_slot})---so a decay into the quark or lepton
sector cannot proceed by continuous relaxation but requires a barriered
reconnection (the $T_{2,2}\!\to\!T_{2,3}$ threading wall of
Paper~XVIII~\cite{Paper18}); the relic is therefore at least
\emph{metastable}. Its own $\Ztwo$ would render it \emph{absolutely}
stable---the standard dark-parity architecture---but only if that $\Ztwo$
is a conserved dynamical symmetry rather than merely an outer
automorphism, which we do not establish here. Stability is thus firm
against unwinding and metastable-by-barrier against decay, with absolute
stability conditional on the $\Ztwo$ being dynamical; only \emph{existence}
is numerically open.
\item \textbf{The branching ratio $r$.} Taken $\approx1$ from comparable
subgroup orders; not derived. The bookkeeping of
Paper~XVIII~\cite{Paper18} shows a dihedral object is \emph{not} a
\emph{forced} per-baryon byproduct (the pion saturates the Hopf-charge and
component books), so $r$ must come from the collective freeze-out
branching, not a forced-byproduct count.
\item \textbf{The annihilation efficiency $f$ and the coldness threshold}
on $\cdir$ (Section~\ref{P21:sec:dm}), both taken order unity/plausible.
\end{enumerate}
Until (i) is settled, the abundance is a scenario, not a result, and the
recombination cold-dark-matter budget remains an input
(Paper~VII~\cite{Paper7}). Results~\ref{P21:thm:noE7} and the sterile-Dirac
character of the D-type sector stand independently of all of the above.

%══════════════════════════════════════════════════════════════════════
\appendix
\section{T-matrix closure engine}
\label{P21:app:engine}
%══════════════════════════════════════════════════════════════════════
The closure computation (Proposition~\ref{P21:prop:noclosure}) is performed
by an exact rational engine
(\path{src_paper21/dtype_cft_tower_engine.py}) that first validates
against the three Standard-Model sectors
($(E_6)_1\to3$, $\SU(2)_1\to6$, $\SU(2)_3\to10$) and then reports no
closing integer for $\SO(10)_1$, $\SO(14)_1$, and all even-level
$\SU(2)_k$. The $a^{-3}$ freeze test (Section~\ref{P21:sec:dm}) uses the
inertial director-network simulation
(\path{src_paper7/dm_polarization/defect_network_sim_3d.py}); the
abundance arithmetic~\eqref{P21:eq:adm} is
\path{src_paper21/adm_abundance.py}.

%══════════════════════════════════════════════════════════════════════
\begin{thebibliography}{99}

\bibitem{Paper1}
F.~Manfredi,
\textit{The Density-Feedback Faddeev--Niemi Hopfion:
  Exact Algebraic Predictions for Standard-Model Parameters},
Zenodo (2026).
\href{https://doi.org/10.5281/zenodo.19342027}{doi:10.5281/zenodo.19342027}

\bibitem{Paper5}
F.~Manfredi,
\textit{Colour Confinement, the QCD Scale, and Hadronic Structure
  from the Density-Feedback Hopfion},
Zenodo (2026).
\href{https://doi.org/10.5281/zenodo.19638857}{doi:10.5281/zenodo.19638857}

\bibitem{Paper7}
F.~Manfredi,
\textit{Emergent Gravity, Inflation, Dark Energy, and the Weak
  Equivalence Principle from the Density-Feedback Hopfion Condensate},
Zenodo (2026).
\href{https://doi.org/10.5281/zenodo.19646953}{doi:10.5281/zenodo.19646953}

\bibitem{Paper15}
F.~Manfredi,
\textit{The $Q_H=3$ Sector of the Density-Feedback Hopfion},
Zenodo (2026).
\href{https://doi.org/10.5281/zenodo.20691001}{doi:10.5281/zenodo.20691001}

\bibitem{Paper16}
F.~Manfredi,
\textit{Quark Generation Masses in the Density-Feedback Hopfion:
A Survey of Ruled-Out Mechanisms and the $Q_H=3$ Gradient-Flow Landscape},
Preprint (2026).
\href{https://doi.org/10.5281/zenodo.21012650}{doi:10.5281/zenodo.21012650}

\bibitem{Paper17}
F.~Manfredi,
\textit{The $Q_H=1$ Sector of the Density-Feedback Hopfion:
Topology, Group Structure, Colour Exclusion},
Zenodo (2026).
\href{https://doi.org/10.5281/zenodo.21013412}{doi:10.5281/zenodo.21013412}

\bibitem{Paper18}
F.~Manfredi,
\textit{Preimage Topology and Confinement: The $Q_H=3$ Sector of the
Density-Feedback Hopfion},
Zenodo (2026).
\href{https://doi.org/10.5281/zenodo.21047750}{doi:10.5281/zenodo.21047750}

\bibitem{Paper19}
F.~Manfredi,
\textit{Dynamical Origin of Quark Masses in the Density-Feedback Hopfion},
Zenodo (2026).
\href{https://doi.org/10.5281/zenodo.21225452}{doi:10.5281/zenodo.21225452}

\bibitem{Paper20}
F.~Manfredi,
\textit{The Dynamical Layer of the Quark Sector: $E_6$-Native
Generations, the Isospin Mass-Scale, and Static Residuals as Strong
Dynamics},
Zenodo (2026).
\href{https://doi.org/10.5281/zenodo.23047960}{doi:10.5281/zenodo.23047960}

\bibitem{Paper22}
F.~Manfredi,
\textit{Dark Genesis: Co-genesis of Baryons and a Sterile Dirac Dark Sector
at the First-Order Confinement Transition of the Density-Feedback Hopfion},
Zenodo (2026) --- DRAFT.
\href{https://doi.org/10.5281/zenodo.23319779}{doi:10.5281/zenodo.23319779}

\bibitem{CIZ1987}
A.~Cappelli, C.~Itzykson and J.-B.~Zuber,
\textit{The A-D-E classification of minimal and $A_1^{(1)}$ conformal
  invariant theories},
Commun.\ Math.\ Phys.\ \textbf{113}, 1--26 (1987).
\href{https://doi.org/10.1007/BF01221394}{doi:10.1007/BF01221394}

\end{thebibliography}

\end{document}
